\section{Step 48 Clip the transform result, then publish}\label{h263-transform-s08}
Clip forward coefficients to the exercise's twelve-bit range, or inverse output
to its signed nine-bit range. Store the sixty-four signed halfwords privately,
then copy the completed 128-byte output to the caller. Return successful count
64 only after publication. Keep every caller byte unchanged on rejection.

Do not clip inverse output to unsigned picture bytes here. INTER reconstruction
still needs to add prediction to the signed inverse residual before final sample
clipping. Quantization and reconstruction will make that next boundary explicit.

\paragraph{Observe} Inverse clipping can return negative values down to $-256$.
The flat inverse observation returns 128 at every position.
\Needspace{8\baselineskip}
\paragraph{Check} Why must a negative inverse residual survive this function?
\begin{itemize}
\item[A.] It may need to reduce the later predicted sample.
\item[B.] It is an output pointer.
\item[C.] H.263 stores every reconstructed picture sample as a negative number.
\end{itemize}
